x

\subsubsection{Fulkerson\textless T\textgreater: алгоритм Фалкерсона}

\textbf{Вход:} \texttt{Graph<T> const\&} (конструктор).

\textbf{Выход:} \texttt{vector<int>}~-- новые номера вершин после
топологической сортировки.

\textbf{Описание:} На каждой итерации метод \texttt{findRoots} смотрит
матрицу смежности и находит вершины без входящих ребер, которые не были удалены.
Каждой такой вершине присваивается очередной номер, затем она помечается
удалённой. Шаги повторяются, пока не упорядочены все вершины.

Если вершин без входящих дуг нет, а вершины ещё остались~-- значит в графе
есть цикл и сортировка невозможна.

\begin{lstlisting}[style=cstyle, caption={Поиск корней и вычисление порядка (fulkerson.cpp)}]
template <typename T>
vector<int> Fulkerson<T>::findRoots(
    const vector<bool>& removed) const
{
    auto adj = graph.getAdjMatrix();
    vector<int> roots;
    for (int v = 0; v < n; v++) {
        if (removed[v]) continue;
        bool hasIncoming = false;
        for (int from = 0; from < n; from++) {
            iterCount++;
            if (!removed[from] && from != v
                && adj[from][v] != 0) {
                hasIncoming = true;
                break;
            }
        }
        if (!hasIncoming) roots.push_back(v);
    }
    return roots;
}

template <typename T>
vector<int> Fulkerson<T>::computeOrder() const {
    iterCount = 0;
    vector<bool> removed(n, false);
    vector<int> order(n, -1);
    int counter = 0;
    while (counter < n) {
        vector<int> roots = findRoots(removed);
        if (roots.empty()) { /* цикл в графе */ return {}; }
        for (int v : roots) {
            order[v] = counter++;
            removed[v] = true;
        }
    }
    return order;
}
\end{lstlisting}

После сортировки вызывается \texttt{Graph<T>::reorder(order)},
и выводятся перенумерованные матрицы смежности и весов.

\subsubsection{FloydWarshall\textless T\textgreater: алгоритм Флойда-Уоршалла}

\textbf{Вход:} \texttt{Graph<T> const\&} (конструктор).

\textbf{Выход:} матрица расстояний \texttt{distMatrix} и матрица
следующих вершин \texttt{nextMatrix} для восстановления пути.

\textbf{Описание:} Инициализация:
$d[i][i] = 0$, $d[i][j] = w(i,j)$ если ребро есть, иначе $+\infty$;
$\texttt{next}[i][j] = j$ при наличии ребра, $-1$~-- иначе.

Главный тройной цикл перебирает промежуточные вершины $k$:

\begin{lstlisting}[style=cstyle, caption={Основной цикл Флойда-Уоршалла (warshall.cpp)}]
template <typename T>
void FloydWarshall<T>::compute() {
    const T INF = getPosINF();
    iterCount = 0;
    for (int k = 0; k < n; k++)
        for (int i = 0; i < n; i++)
            for (int j = 0; j < n; j++) {
                iterCount++;
                if (distMatrix[i][k] == INF
                    || distMatrix[k][j] == INF) continue;
                T newDist = distMatrix[i][k] + distMatrix[k][j];
                if (newDist < distMatrix[i][j]) {
                    distMatrix[i][j] = newDist;
                    nextMatrix[i][j] = nextMatrix[i][k];
                }
            }
}
\end{lstlisting}

Путь восстанавливается методом \texttt{getPath(from, to)}: идём по
\texttt{nextMatrix} от \texttt{from} до \texttt{to}.

\subsubsection{Сравнение алгоритмов по числу итераций}

Пункт меню~11 запускает Фалкерсона и Флойда-Уоршалла на одном графе
и сравнивает счётчики \texttt{iterCount}~-- каждый из них растёт
при каждом просмотре ячейки матрицы. Так можно честно сравнить
трудоёмкость без привязки к железу.
